% Part: first-order-logic
% Chapter: lindstrom
% Section: lindstrom-proof

\documentclass[../../../include/open-logic-section]{subfiles}

\begin{document}

\olfileid{mod}{lin}{prf}

\olsection{Teorema Lindstr\"om}

\begin{lem}
\ollabel{lem:lindstrom}
Upamana $!E \in L(\Lang{L})$, kanthi $\Lang{L}$ winates, lan
upamana uga ana $n \in \Nat$ saengga kanggo sembarang rong
!!{structure}s $\Struct{M}$ lan~$\Struct{N}$, yen $\Struct{M} \equiv_n
\Struct{N}$ lan $\Struct{M} \models_L !E$, mula $\Struct{N}
\models_L !E$. Banjur $!E$ ekuivalen karo !!{sentence}
orde kapisan, yaiku ana $!D$ orde kapisan saengga
$\Mod(L){!E} = \Mod(L){!D}$.
\end{lem}

\begin{proof}
Upamana $n$ saengga sembarang rong !!{structure}s kang ekuivalen-$n$,
$\Struct{M}$ lan $\Struct{N}$, sarujuk marang biji bebener~$!E$.
Elinga \olref[bas][pis]{prop:qr-finite}: nganti ekuivalensi logis,
mung ana cacah winates !!{sentence}s orde kapisan ing
!!{language} winates kang pangkat kuantoré ora ngluwihi~$n$.
Saiki, kanggo
saben !!{structure}~$\Struct{M}$ kang tetep, upamana
$!D_{\Struct{M}}$ minangka konjungsi kabeh !!{sentence}s
orde kapisan~$!E$ kang bener ing~$\Struct{M}$ lan
$\QuantRank{!E} \le n$ (konjungsi iki winates), saengga
$\Struct{N} \models !D_{\Struct{M}}$ yen lan mung yen
$\Struct{N} \equiv_n \Struct{M}$. Banjur tetapna
$!D = \textstyle\bigvee
\Setabs{!D_{\Struct{M}}}{\Struct{M} \models_L !E}$; disjungsi iki
uga winates (nganti ekuivalensi logis).

Saka kene, kasimpulan $\Mod(L){!E} = \Mod(L){!D}$ tumut.
Pancen, yen $\Struct{N} \models_L !D$, mula kanggo sawijining
$\Struct{M} \models_L !E$ kita nduweni
$\Struct{N} \models !D_{\Struct{M}}$, mula uga
$\Struct{N} \models_L !E$ (miturut hipotesis lema).
Kosokbaline, yen $\Struct{N} \models_L !E$, mula
$!D_\Struct{N}$ dadi salah siji suku disjungsi ing $!D$,
lan amarga $\Struct{N} \models !D_\Struct{N}$, uga
$\Struct{N} \models_L !D$.
\end{proof}

\begin{thm}[Teorema Lindstr\"om]
  \ollabel{thm:lindstrom} Upamana $\tuple{L, \models_L}$ nduweni
  Sipat Kekompakan lan Sipat L\"owenheim--Skolem. Banjur
  $\tuple{L, \models_L} \le \tuple{F, \models}$ (mula
  $\tuple{L, \models_L}$ ekuivalen karo logika orde kapisan).
\end{thm}

\begin{proof}
Miturut \olref{lem:lindstrom}, cukup dibuktekake yen kanggo sembarang
$!E \in L(\Lang{L})$, kanthi $\Lang{L}$ winates, ana $n \in \Nat$
saengga kanggo sembarang rong !!{structure}s $\Struct{M}$
lan~$\Struct{N}$: yen $\Struct{M} \equiv_n \Struct{N}$, mula
$\Struct{M}$ lan $\Struct{N}$ sarujuk marang~$!E$. Amarga iku,
$!E$ banjur ekuivalen karo !!{sentence} orde kapisan, lan saka
kene $\tuple{L, \models_L} \le \tuple{F, \models}$ tumut.
Amarga kita makarya ing !!{language} relasional murni kang winates,
miturut \olref[bas][pis]{thm:b-n-f}, pratelan
$\Struct{M} \equiv_n \Struct{N}$ bisa diganti nganggo pratelan
aljabar kang cocog, yaiku $I_n(\emptyset,\emptyset)$.

Kanggo $!E$ kang diwenehake, upamana, kanggo nggoleki cengkah,
yen kanggo saben $n$ ana !!{structure}s $\Struct{M}_n$ lan
$\Struct{N}_n$ saengga $I_n(\emptyset, \emptyset)$,
nanging (umpamane) $\Struct{M}_n \models_L !E$ dene
$\Struct{N}_n \not\models_L !E$. Miturut Sipat Isomorfisme
kita bisa nganggep yen kabeh $\Struct{M}_n$ napsirake konstanta
basa kasebut nganggo objek kang padha; saliyane iku, amarga
ing basa kasebut mung ana cacah winates !!{sentence}s atomik,
kita uga bisa nganggep yen kabeh mau nyukupi !!{sentence}s
atomik kang padha (yen ora, kita bisa njupuk subsekuens saka
$\Struct{M}$). Upamana $\Struct{M}$ minangka gabungan kabeh
$\Struct{M}_n$, yaiku !!{structure} minimal tunggal kang
nduweni saben $\Struct{M}_n$ minangka substruktur. Kaya ing
bukti \olref[lsp]{thm:abstract-p-isom}, upamana
$\Struct{M}^*$ minangka perluasan saka $\Struct{M}$ kanthi
!!{domain} $\Domain{M} \cup
\Domain{M}^{<\omega}$, ing !!{language} kang diperluas
kanthi predikat panyambungan $P$ lan~$Q$.

Kanthi cara kang padha, tetapna $\Struct{N}_n$, $\Struct{N}$
lan $\Struct{N}^*$. Saiki upamana $\Struct{M}$ minangka
!!{structure} kang !!{domain}-e ngemot !!{domain}s saka
$\Struct{M}^*$ lan $\Struct{N}^*$ uga wilangan asli~$\Nat$
kalebu urutan alamiahe~$\le$, ing !!{language} kanthi predikat
tambahan kang nglambangake !!{domain}s
$\Domain{M}$, $\Domain{N}$, $\Domain{M}^{<\omega}$ lan
$\Domain{N}^{<\omega}$ uga predikat kang nyandi ranah saka
$\Struct{M}_n$ lan $\Struct{N}_n$ ing pangerten iki:
\begin{align*}
  \Domain{M_n} & = \Setabs{a \in \Domain{M}}{R(a, n)}; &
  \Domain{N_n} & = \Setabs{a \in \Domain{N}}{S(a,n)}; \\
  \Domain{M}^{<\omega}_n & = \Setabs{a \in \Domain{M}^{<\omega}}{R(a,n)}; &
  \Domain{N}^{<\omega}_n & = \Setabs{a \in \Domain{N}^{<\omega}}{S(a,n)}.
\end{align*}
!!{structure}~$\Struct{M}$ iki uga nduweni relasi telung panggonan
$J$ saengga $J(n, \mathbf{a}, \mathbf{b})$ lumaku yen lan mung
yen $I_n(\mathbf{a}, \mathbf{b})$.

Saiki ana !!a{sentence}~$!D$ ing !!{language}~$\Lang{L}$ kang
ditambahi $R$, $S$, $J$, lan sapiturute, kang nyatakake yen
$\le$ iku urutan linear diskret kang nduweni unsur kapisan
nanging ora nduweni unsur pungkasan, lan yen
$\Struct{M}_n \models !E$, $\Struct{N}_n \not\models !E$,
lan kanggo saben $n$ ing urutan kasebut,
$J(n, \mathbf{a}, \mathbf{b})$ lumaku yen lan mung yen
$I_n(\mathbf{a}, \mathbf{b})$.

Nganggo Sipat Kekompakan, kita bisa nemokake model
$\Struct{M}^*$ saka $!D$ kang urutane ngemot unsur
nonstandar~$n^*$. Mula, mligine, $\Struct{M^*}$ ngemot
sub!!{structure}s $\Struct{M_{n^*}}$ lan
$\Struct{N_{n^*}}$ saengga $\Struct{M_{n^*}}
\models_L !E$ lan $\Struct{N_{n^*}} \not\models_L !E$.
Nanging saiki kita bisa nemtokake himpunan $\mathcal{I}$
saka pasangan tupel-$k$ saka $\Domain{M_{n^*}}$ lan
$\Domain{N_{n^*}}$ kanthi netepake
$\tuple{\mathbf{a}, \mathbf{b}} \in \mathcal{I}$ yen lan mung
yen $J(n^*-k, \mathbf{a}, \mathbf{b})$, ing ngendi $k$
iku dawané $\mathbf{a}$ lan $\mathbf{b}$. Amarga $n^*$
nonstandar, kanggo saben $k$ standar kita nduweni
$n^* - k >0$, lan himpunan $\mathcal{I}$ iki dadi seksi
yen $\Struct{M_{n^*}} \simeq_p
\Struct{N_{n^*}}$. Nanging miturut
\olref[lsp]{thm:abstract-p-isom}, $\Struct{M_{n^*}}$
iku ekuivalen-$L$ karo $\Struct{N_{n^*}}$, lan iki cengkah.
\end{proof}

\end{document}
